\documentclass[11pt,a4paper]{article}
\usepackage[margin=25mm]{geometry}
\usepackage[T1]{fontenc}
\usepackage{hyperref}
\newcommand{\code}[1]{\texttt{\detokenize{#1}}}
\title{Blame: Architecture Summary}
\author{Blame Engineering}
\date{October 7, 2026}
\begin{document}
\maketitle

\section{System boundaries}
Blame combines a React application, an Electron desktop host, an AdonisJS
backend, and a Cline SDK adapter. The application renders backlog, conversation,
review, and preview state. It accesses desktop capabilities through the platform
interface and the preload bridge. Electron IPC registration is assembled by
\code{packages/desktop/src/main/ipc.ts}; workspace handlers delegate Git
operations to the workspace service. Renderer code must not introduce direct
Electron or host filesystem dependencies.

\code{packages/agent-runtime} adapts SDK events and coordinates agent sessions,
approvals, tools, cancellation, and delivery. SDK-specific configuration and
payload normalization stay in that adapter layer. \code{packages/harness-core}
contains framework-independent harness contracts and reference authorities.
\code{packages/preview-contract} defines shared preview contracts, while
\code{packages/preview-runtime} implements lifecycle coordination and provider
adapters. The backend owns authentication, persisted project and conversation
data, encrypted provider credentials, and hosted integrations.

Dependencies should point from UI and infrastructure towards application
contracts and domain rules. Transport handlers translate requests; services
coordinate side effects; domain rules determine outcomes. Structure should
follow an actual responsibility rather than a file's length.

\section{Current Light Semantic Harness}
\subsection{Model interface and runtime ownership}
The configured Light path exposes five core primitives:
\code{search_code}, \code{read_code}, \code{edit_code},
\code{run_task}, and \code{finish_task}. Preview and other configured tools
may extend a session. The runtime owns target resolution, observations,
source revisions, guarded effects, and execution. Local operations use
\code{LightWorkspace}; managed operations use a shared executor port.

Syntax adapters supply declarations and byte ranges. Reference search currently
uses lexical fallback; this is not a compiler-level reference graph.
Aliases bind observed targets to content hashes. Changed target bodies require
fresh observation, while unchanged declarations can relocate after neighboring
changes. Read coverage is scoped to content and context generation.

\subsection{Results, edits, and recovery}
The production tool envelope is \code{LightToolResult}, with
\code{ok}, \code{blocked}, and \code{failed} outcomes. Command observations
must be read separately from transport success. A journal persists changes,
validations, aliases, coverage, and terminal state.

\code{edit_code} prepares and validates a batch before revision and preimage
checks. A cooperating lock, durable pending record, atomic file updates, and
owned rollback support interrupted-write recovery. Recovery preserves detected
newer external changes. These controls do not establish global filesystem
isolation against programs that ignore the lock.

\subsection{Validation and completion}
Completion differs from SDK termination. The finish service considers current
effects and relevant checks already requested in the session. Fresh code-check
failures remain blocking; missing admitted validation yields unverified delivery.
It does not invent an unrelated build. Test execution, selected scope, freshness,
and admitted verification are distinct: observed passing tests currently retain
unknown validation status when selection and freshness are not established.

The inspected default task contract permits no source effect. Free-text
acceptance criteria are not mechanically solved by the finish service.
Verified delivery means coverage under the current validation policy, not
proof of arbitrary functional requirements.

\subsection{Historical reference contracts}
The core package retains authority, result, transaction, and evidence-planner
reference contracts from earlier designs. Execution lifecycle functionality
is integrated, and completion contracts are shared as types. The Light result
envelope and finish path do not invoke the core evidence planner or instantiate
the archived four-authority protocol. Exported reference implementations must
not be presented as proof of production wiring.

The detailed analysis, evidence limits, benchmark results, and recommendations
are in \code{docs/semantic-harness-study.tex}.

\section{Workspace and preview execution}
Local execution uses native worktrees on Linux, Lima on macOS, and WSL on
Windows. Cloud-preview work keeps the agent and worktree native and explicitly
synchronizes a snapshot to a Vercel Sandbox. Transport or browser readiness is
insufficient to prove that the application is ready for the current source and
environment revisions. Preview contracts and lifecycle services remain the
boundary for provider-specific behavior.

Deploying to a user's Vercel account is a separate integration from managed
sandbox previews. The backend handles account authorization and protected
credentials; deployment snapshot and preview access rules remain explicit.
These integrations must not become assumptions of the portable domain layer.

\section{Persistence and presentation}
Streaming projections update the UI immediately. Backlog and quick-chat writes
are coalesced per target through a shared delayed queue. Terminal writes retain
their durable outbox, finalization, acknowledgment, and replay paths. A queue is
a timing mechanism; it does not replace persistence ownership or retry policy.

Markdown editing and preview rendering have separate modules. Native browser
sessions retain a shared registry so warm-up and visible panels can reuse a
session. The registry controls reference counts, creation, hiding, retention,
and destruction. Conversation presentation interprets tool details before
rendering them; presentation must not redefine runtime evidence or completion.

\section{Documentation scope}
This summary replaces historical architecture drafts, including the
decommissioned Semantic Workspace Harness proposal. The current production
direction is the Light Semantic Harness. Historical benchmark traces are kept
in Git history and are not treated as new live-provider measurements.
Source code and executable tests determine implemented behavior.
\end{document}
